\documentclass[10pt,a4paper]{amsart}
\usepackage[margin=19mm]{geometry}
\usepackage{amsmath,amssymb,mathtools}
\usepackage[scr=rsfs,cal=euler]{mathalfa}
\usepackage{xfrac}
\usepackage[unicode,hidelinks,bookmarks=false]{hyperref}
\newcommand{\ie}{i.e.\ }
\newcommand{\cf}{cf.\ }
\newcommand{\resp}{respectively\ }
\newcommand{\Subsection}{\subsection}
\providecommand{\nobreakdash}{\nobreak\hbox{-}}
\setlength{\emergencystretch}{2em}
\multlinegap0pt
\usepackage{float}
\usepackage{amssymb}
\usepackage{dsfont}
\usepackage{xcolor}
\usepackage{tikz}
\usepackage{mathtools}
\usepackage{bm}
\newtheorem{theorem}{Theorem}
\newcommand{\N}{\mathbb{N}}
\newcommand{\Om}{\Omega}
\newcommand{\dsp}{\displaystyle}
\newcommand{\mH}{\mrm{H}}
\newcommand{\mL}{\mrm{L}}
\newcommand{\mrm}[1]{\mathrm{#1}}
\title{A few techniques to achieve invisibility in waveguides}
\author{Lucas Chesnel}
\date{Source: arXiv:2511.09172v2 (CC BY 4.0)}
\begin{document}
\maketitle

\noindent\emph{Source and licence.} A few techniques to achieve invisibility in waveguides. Version arXiv:2511.09172v2 (CC BY 4.0), distributed by arXiv under the Creative Commons Attribution 4.0 International licence, http://creativecommons.org/licenses/by/4.0/, as recorded in the arXiv licence metadata for this exact version and shown on the abstract page https://arxiv.org/abs/2511.09172v2. Full licence notice: \texttt{licenses/arxiv-2511.09172-CC-BY-4.0.txt}.\par\smallskip

\noindent\textit{Editorial scope.} Counted unit: the article's theorem theorem (source0.tex lines 1093--1099). The article's complete original proof of it is delivered with it (source0.tex lines 1100--1169, one \texttt{proof} environment of 70 lines). No mathematics is written, repaired or re-derived: the statement and the proof are the article's own lines, in the article's own notation, and no retained line is substituted or re-flowed.  Retained uncounted as the source-local context the proof needs: lines 928--934; lines 954--960; lines 987--989; lines 991--997.  Nothing was omitted from any retained span, so the package carries no omission record.  No retained line contains a citation, so the wrapper carries no bibliography and no bibliography entry.  No retained line contains a figure environment, an image-inclusion command or drawing code, so no image is delivered and the package declares no asset dependency.  Editorial formatting only, in addition: an AMS \texttt{amsart} preamble that restates the article's own declarations and macros used by the retained lines, namely the environments \texttt{theorem} on separate counters, together with the standard packages those lines need.  The retained environments print in this document as Theorem 1 in order of appearance, whereas the article prints the same items with the numbering of its own full document.  The complete native source of the article as distributed in the arXiv e-print for this exact version is bundled unchanged as \texttt{sources/188956/source0.tex}.
\begin{equation}\label{WaveguidePbDissipationBounded}
\begin{array}{|rcll}
\multicolumn{4}{|l}{\mbox{Find }u_\eta\in\mH^1_0(\Om_L;\Gamma)\mbox{ such that }}\\[2pt]
-\Delta u_\eta-(k^2+ik\eta)u_\eta&=&f&\mbox{ in }\Om_L \\[2pt]
\cfrac{\partial u_\eta}{\partial \nu} & = & \Lambda_\pm^\eta(u_\eta|_{\Sigma_{\pm L}}) &\mbox{ on }\Sigma_{\pm L}
\end{array}
\end{equation} 

\begin{equation}\label{WaveguidePbBounded}
\begin{array}{|rcll}
\multicolumn{4}{|l}{\mbox{Find }u\in\mH^1_0(\Om_L;\Gamma)\mbox{ such that }}\\[2pt]
-\Delta u-k^2u&=&f&\mbox{ in }\Om_L \\[2pt]
\cfrac{\partial u}{\partial \nu} & = & \Lambda_\pm(u|_{\Sigma_{\pm L}}) &\mbox{ on }\Sigma_{\pm L}
\end{array}
\end{equation} 

\begin{equation}\label{DefOpAkOut}
(A^{\mrm{out}}(k)u,v)_{\mH^1(\Om_L)} =a^{\mrm{out}}(u,v),\qquad\forall u,v\in \mH^1_0(\Om_L;\Gamma).
\end{equation}

\begin{theorem}\label{ThmFredholm}
For $k\in(\pi;+\infty)\setminus\{\N\pi\}$, the operator $A^{\mrm{out}}(k)$ decomposes as 
\[
A^{\mrm{out}}(k)=B^{\mrm{out}}+K
\]
where $B^{\mrm{out}}:\mH^1_0(\Om_L;\Gamma)\to\mH^1_0(\Om_L;\Gamma)$ is an isomorphism and $K:\mH^1_0(\Om_L;\Gamma)\to\mH^1_0(\Om_L;\Gamma)$ is compact ($B^{\mrm{out}}$ and $K$ are allowed to depend on $k$).
\end{theorem}

\begin{theorem}
Fix $k\in(\pi;+\infty)\setminus\{\N\pi\}$ and assume that the operator $A^{\mrm{out}}(k)$ defined in (\ref{DefOpAkOut}) is injective. Then there is $\eta_0>0$ such that we have
\[
\|u-u_\eta\|_{\mH^1(\Om_L)} \le C\eta\,\|f\|_{\mL^2(\Om_L)},\qquad \forall \eta\in(0;\eta_0],
\]
with a constant $C>0$ which is independent of $\eta$. Here $u$, $u_\eta$ are the functions solving respectively (\ref{WaveguidePbBounded}), (\ref{WaveguidePbDissipationBounded}).
\end{theorem}

\begin{proof}
Let $F$ be the element of $\mH^1_0(\Om_L;\Gamma)$ such that 
\[
(F,v)_{\mH^1(\Om_L)}=\ell(v),\qquad\forall v\in\mH^1_0(\Om_L;\Gamma).
\]
Denote also by $A^\eta:\mH^1_0(\Om_L;\Gamma)\to\mH^1_0(\Om_L;\Gamma)$ the operator such that
\[
(A^\eta v,w)_{\mH^1(\Om_L)}=\dsp\int_{\Om_L}\hspace{-0.2cm}\nabla v\cdot\nabla\overline{w}-(k^2+ik\eta) v\overline{w}\,dxdy-\langle\Lambda^\eta_+(v|_{\Sigma_{L}}),\overline{w} \rangle_{\Sigma_L}-\langle\Lambda^\eta_-(v|_{\Sigma_{-L}}),\overline{w} \rangle_{\Sigma_{-L}}
\]  
for all $v,w\in\mH^1_0(\Om_L;\Gamma)$. We have, for all $\eta>0$,
\[
A^\eta u_\eta=F=A^{\mrm{out}} u. 
\]
This gives $A^{\mrm{out}} (u-u_\eta)=(A^\eta-A^{\mrm{out}})u_\eta$. Since $A^{\mrm{out}}$ is assumed to be injective, Theorem \ref{ThmFredholm} together with the Fredholm alternative guarantee that $A^{\mrm{out}}$ is an isomorphism of $\mH^1_0(\Om_L;\Gamma)$. Thus we can write
\begin{equation}\label{Relation1}
u-u_\eta=(A^{\mrm{out}})^{-1} (A^\eta-A^{\mrm{out}})u_\eta.
\end{equation}
Let us estimate
\begin{equation}\label{DefNormDif}
\|A^\eta-A^{\mrm{out}}\| =\underset{v\in\mH^1_0(\Om_L;\Gamma)\setminus\{0\}}{\sup}\,\cfrac{\|(A^\eta -A^{\mrm{out}}) v\|_{\mH^1(\Om_L)}}{\|v\|_{\mH^1(\Om_L)}}\,.
\end{equation}
From the definition of $A^{\mrm{out}}$, $A^\eta$, one finds, for $v,w\in\mH^1_0(\Om_L;\Gamma)$, 
\[
\begin{array}{ll}
&((A^\eta -A^{\mrm{out}}) v,w)_{\mH^1(\Om_L)} \\[8pt]
=&-ik\eta\dsp\int_{\Om_L} v\,\overline{w}\,dxdy -\langle\Lambda^\eta_+(v|_{\Sigma_{L}})-\Lambda_+(v|_{\Sigma_{L}}),\overline{w} \rangle_{\Sigma_L}-\langle\Lambda^\eta_-(v|_{\Sigma_{-L}})-\Lambda_-(v|_{\Sigma_{-L}}),\overline{w }\rangle_{\Sigma_{-L}}. 
\end{array}
\]
Taking $w=(A^\eta -A^{\mrm{out}}) v$ in this equality and using the continuity of the trace mappings from $\mH^1(\Om_L)$ to $\mH^{1/2}(\Sigma_{\pm L})$, we deduce the existence of a constant
$C>0$, which may change from one line to another below but which remains independent of $\eta\in(0;\eta_0]$, such that
\begin{equation}\label{DefNormDif2}
\begin{array}{l}
\|(A^\eta -A^{\mrm{out}}) v\|_{\mH^1(\Om_L)} \le k\eta\,\|v\|_{\mH^1(\Om_L)}\\[9pt]
\hspace{1.6cm}+C\,(\|\Lambda^\eta_+(v|_{\Sigma_{L}})-\Lambda_+(v|_{\Sigma_{L}})\|_{\mH^{-1/2}(\Sigma_{L})}+\|\Lambda^\eta_-(v|_{\Sigma_{-L}})-\Lambda_-(v|_{\Sigma_{-L}})\|_{\mH^{-1/2}(\Sigma_{-L})}).
\end{array}
\end{equation} 
Let us estimate the second term in the right-hand side above. For $\varphi,\varphi'\in \mH^{1/2}_{00}(\Sigma_{\pm L})$, we can write
\begin{equation}\label{EstimaDtNCont0Bis}
|\langle \Lambda_\pm^\eta(\varphi)-\Lambda_\pm(\varphi),\overline{\varphi'}\rangle_{\Sigma_{\pm L}}| \le \dsp\sum_{n=1}^{+\infty} |\beta_n^\eta-\beta_n|\,|(\varphi,\varphi_n)_{\mL^2(\Sigma_{\pm L})}|\,|(\varphi',\varphi_n)_{\mL^2(\Sigma_{\pm L})}|. 
\end{equation}
Working from the expressions of $\beta_n^\eta$, $\beta_n$, one shows that for any given $\eta_0>0$, there is a constant $C>0$ independent of $\eta\in(0;\eta_0]$ such that 
\[
\underset{n\in\N^\ast}{\sup} \cfrac{|\beta_n^\eta-\beta_n|}{n\pi} \le C\eta.
\]
Exploiting this in (\ref{EstimaDtNCont0Bis}) gives
\[
\begin{array}{rcl}
|\langle \Lambda_\pm^\eta(\varphi)-\Lambda_\pm(\varphi),\overline{\varphi'}\rangle_{\Sigma_{\pm L}}|
&\le &  \dsp C\eta\,\left(\dsp\sum_{n=1}^{+\infty}n\pi|(\varphi,\varphi_n)_{\mL^2(\Sigma_{\pm L})}|^2\right)^{1/2} \left(\dsp\sum_{n=1}^{+\infty}n\pi|(\varphi',\varphi_n)_{\mL^2(\Sigma_{\pm L})}|^2\right)^{1/2} \\[12pt]
&\le &  \dsp C\eta\,\|\varphi\|_{\mH^{1/2}(\Sigma_{ L})}\|\varphi'\|_{\mH^{1/2}(\Sigma_{ L})}.
\end{array}
\]
From this, we deduce the estimates
\begin{equation}\label{DefNormDif3}
\|\Lambda^\eta_\pm(v|_{\Sigma_{\pm L}})-\Lambda_\pm(v|_{\Sigma_{\pm L}})\|_{\mH^{-1/2}(\Sigma_{\pm L})}\le C\,\eta\,\|v\|_{\mH^{1/2}(\Sigma_{\pm L})} \le C\,\eta\,\|v\|_{\mH^{1}(\Om_L)}.
\end{equation}
By combining (\ref{DefNormDif2}), (\ref{DefNormDif3}) and  (\ref{DefNormDif}), one gets, for $\eta\in(0;\eta_0]$,
\begin{equation}\label{Relation2}
\|A^\eta-A^{\mrm{out}}\| \le C \eta,
\end{equation}
Then, gathering (\ref{Relation2}) and (\ref{Relation1}), we find
\begin{equation}\label{FirstEstimate}
\|u-u_\eta\|_{\mH^1(\Om_L)} \le C \eta\, \|u_\eta\|_{\mH^1(\Om_L)}.
\end{equation}
Finally, by using the inequality $\|u_\eta\|_{\mH^1(\Om_L)}\le \|u-u_\eta\|_{\mH^1(\Om_L)}+\|u\|_{\mH^1(\Om_L)}  $ in (\ref{FirstEstimate}), we obtain, for $\eta$ small,
\[
\|u-u_\eta\|_{\mH^1(\Om_L)} \le C\eta\,\|u\|_{\mH^1(\Om_L)}\le C\eta\,\|(A^{\mrm{out}})^{-1}f\|_{\mH^1(\Om_L)} \le C\eta\,\|f\|_{\mL^2(\Om_L)}.
\]
This completes the proof.
\end{proof}

\end{document}
